\chapter{The obstruction constitution}
\label{ch:no-go}

The statements of this chapter are the ones that prevent a forgetful image
from being promoted into its source. Each names what is \emph{not}
determined, and therefore what must be supplied. They are load-bearing:
without them the theorems of the preceding chapters admit false corollaries.

\begin{theorem}[No quantum group from dimension loss]
\label{nogo:dimension-loss}
\StatusProved{} The dimension of the kernel of $V^{\otimes n}\to\Sym^nV$
does not determine an $r$-matrix, an associator, or a quantum group. That
kernel carries no multiplication, no coproduct, no monodromy, and no
coherence equation.
\end{theorem}

\begin{proof}
The kernel is a representation of $\Sigma_n$ and nothing more. Two
algebras with the same underlying graded vector space, hence the same
kernel dimensions in every arity, may have different multiplications and
inequivalent representation categories; the free algebra and any of its
quadratic quotients with the same Hilbert series is an example. A
quantum-group structure is a coproduct together with coherence data, none
of which is a function of a dimension.
\end{proof}

Consequently an averaging or symmetrisation map, which records exactly
this dimension loss, does not produce quantum symmetry. Quantum symmetry
arises only through the exchange data of Chapter~\ref{ch:two-directions}
and a Tannakian reconstruction with its fibre-functor hypotheses stated.

\begin{theorem}[No modular object from topology alone]
\label{nogo:topology}
\StatusProved{} The identity $\partial^2=0$ on
$\overline{\mathcal M}_{g,n}$ assigns no coefficient operations to
boundary strata. A modular master element exists only after cyclic
contractions and clutching-compatible coefficient maps have been supplied.
\end{theorem}

\begin{proof}
$\partial^2=0$ is a statement about the stratification. Given a
coefficient complex $E$ and a degree-one contraction $\rho_F$ for each
codimension-one face, compatible with restriction to corners and
satisfying $d_E\rho_F+\rho_Fd_E=0$, the operator $D=d_E+\sum_F\rho_F$ has
$D^2$ supported on codimension-two strata as the signed sum of incident
composites. Without the maps $\rho_F$ there is no equation to satisfy, and
the maps are not determined by the stratification.
\end{proof}

\begin{theorem}[No bulk reconstruction from a boundary centre]
\label{nogo:bulk}
\StatusProved{} $Z(A)=\RHom_{A^e}(A,A)$ is terminal among algebraic
central actions on the boundary algebra $A$. It does not detect a
decoupled bulk sector. A physical bulk is therefore not determined by the
boundary chiral algebra unless a bulk-to-centre comparison
$\mathcal O^{\mathrm{phys}}_{\mathrm{bulk}}(A)\to Z(A)$ is supplied and
proved to be an equivalence.
\end{theorem}

\begin{proof}
Terminality is the universal property of the centre: any algebra acting
centrally on $A$ factors through $Z(A)$ uniquely up to coherent homotopy.
Terminality is a statement about the boundary category and is insensitive
to any sector that acts trivially on it, so a bulk with trivial boundary
action is invisible to $Z(A)$.
\end{proof}

\begin{theorem}[No classification by central charge]
\label{nogo:central-charge}
\StatusProved{} Central charge determines the Virasoro anomaly
coefficient. It does not determine the chiral operation, the
representation category, the braid monodromy, or the higher-genus
conformal blocks. Distinct holomorphic theories at equal central charge
exist.
\end{theorem}

\begin{theorem}[No identification of anomalies of different types]
\label{nogo:anomaly-types}
\StatusProved{} A current-algebra level, a Virasoro central charge, a
determinant-line curvature, a BRST ghost-number anomaly, and a
Borcherds-product weight are invariants in different theories, living in
different deformation complexes. They may be related by a theorem in a
particular family. They are not one universal scalar, and no identity
among them follows from numerical agreement.
\end{theorem}

\begin{proof}
Each invariant is a class in a distinct complex: an Atiyah-algebra
extension class, a Virasoro cocycle, the curvature of a determinant line,
a class in $H^1$ or $H^2$ of a BRST convolution complex, and a weight in
the character lattice of an automorphic form. A map between two of these
complexes is additional data. In its absence, equality of numbers is a
coincidence of values, not an identification of classes.
\end{proof}

Theorem~\ref{nogo:central-charge} and
Theorem~\ref{nogo:anomaly-types} together retract any programme that
treats a tuple of scalars drawn from these five sources as a single
classifying invariant of a chiral algebra, and any trace identity asserted
among them without a named comparison morphism.

\begin{theorem}[No unnamed comparison of bars]
\label{nogo:unnamed-bar-comparison}
\StatusProved{} An asserted equality $\Barperp(A)=\BarchAss(A)$ has no
meaning without a functor carrying the pointwise multiplication problem to
the chiral-convolution multiplication problem. The two constructions have
different underlying gradings, counting interval subdivisions and
collision trees respectively. Agreement of Euler characteristics is not a
comparison map.
\end{theorem}

\begin{proof}
The two bars are bar constructions over different monoidal products on
different categories, by Notation~\ref{not:three-bars}. A morphism between
them is a monoidal comparison functor together with the datum that it
carries one multiplication to the other; no such functor is canonical, and
by Proposition~\ref{prop:lambda-sector} even a distributive law between
the products exists only on the aligned sector.
\end{proof}

\begin{theorem}[A scalar sequence is not homology data]
\label{nogo:scalar-sequence}
\StatusProved{} A sequence of integers becomes bar homology only after a
chain complex, a differential, a grading convention, and a rank
computation or quasi-isomorphism are specified. Matching initial values,
discriminants, or growth rates constructs no such map.
\end{theorem}

\begin{proof}
Chain-space dimensions and homology dimensions differ by the ranks of the
differentials, which are not functions of the chain dimensions. Explicitly,
Computation~\ref{comp:simple} and Theorem~\ref{thm:exact-audit} exhibit an
algebra whose bar chain dimensions are $3,9,27,81,243$ and whose bar
homology dimensions are $1,0,0,1$; and
Proposition~\ref{prop:one-sequence} exhibits two sequences with the same
discriminant and growth rate that are elementary transforms of one another
and of neither homology.
\end{proof}
